\begin{proof}[Proof of \cref{obj:lem:parametric-floor-infrastructure}]
Fix positive integers \(n,d\) and \(q\in[1/2,1]\).

\begin{enumerate}
\item \emph{Risk calibration.}
Define
\[
a_{\mathrm{risk}}
:=
\min\left\{\frac14,\frac{\sqrt{\log 2}}4\right\},
\qquad
u_{\mathrm{risk}}
:=
\frac{a_{\mathrm{risk}}}{\sqrt n},
\qquad
c_0:=\frac{a_{\mathrm{risk}}^2}{16}.
\]
Since \(\log 2>0\),
\[
0<a_{\mathrm{risk}}\le\frac14,
\qquad
a_{\mathrm{risk}}^2\le\frac{\log 2}{16},
\qquad
c_0>0.
\]
Moreover, \(n\ge1\) implies \(\sqrt n\ge1\) and \((\sqrt n)^2=n\). Consequently,
\[
0<u_{\mathrm{risk}}\le\frac14,
\qquad
n u_{\mathrm{risk}}^2=a_{\mathrm{risk}}^2,
\qquad
16n u_{\mathrm{risk}}^2\le\log 2.
\]
Thus the same positive constant \(a_{\mathrm{risk}}\) meets this budget for every positive sample size.
% lean: CausalSmith.Stat.MarNearcompleteFrontier.parametric_radius_budget

\item \emph{A one-cell pair.}
For each parameter \(\theta\in[0,1]\), define a full-data law \(P_\theta\) by
\[
\begin{gathered}
X=0,\qquad S(0)=S(1)=0,\qquad Y(0)=0,\qquad R=1
\quad\text{almost surely},\\
A\sim\operatorname{Bernoulli}(1/2),\qquad
Y(1)\sim\operatorname{Bernoulli}(\theta),\qquad
A\perp Y(1),\\
(S,Y)=(S(A),Y(A)),
\qquad
Q_\theta:=(P_\theta)_O .
\end{gathered}
\]
The label \(0\) is available because \(d\ge1\). The construction gives randomized treatment independence, balanced assignment, and consistency. Since \(R=1\), outcome arrival is conditionally independent of \(Y\), and every occupied cell has arrival probability one. Hence the requirements in \cref{obj:def:law-class} are satisfied:
\[
P_\theta\in\mathcal M(d,q),
\qquad
\tau(P_\theta)
=\mathbb E_{P_\theta}\{Y(1)-Y(0)\}
=\theta.
\]
The observed law has exactly the following possible nonzero masses:
\[
\begin{aligned}
Q_\theta\{(0,0,0,1,0)\}&=\frac12,\\
Q_\theta\{(0,1,0,1,0)\}&=\frac{1-\theta}{2},\\
Q_\theta\{(0,1,0,1,1)\}&=\frac{\theta}{2},
\end{aligned}
\]
and assigns mass zero to every other record.

Define the risk pair by
\[
P_{\mathrm{risk},0}:=P_{1/2},
\qquad
P_{\mathrm{risk},1}:=P_{1/2+u_{\mathrm{risk}}},
\qquad
Q_{\mathrm{risk},\iota}:=(P_{\mathrm{risk},\iota})_O
\quad(\iota\in\{0,1\}).
\]
Both parameters belong to \([0,1]\), and
\[
\tau(P_{\mathrm{risk},0})=\frac12,
\qquad
\tau(P_{\mathrm{risk},1})=\frac12+u_{\mathrm{risk}},
\qquad
Q_{\mathrm{risk},1}\ll Q_{\mathrm{risk},0}.
\]
For this finite, absolutely continuous pair, the chi-square divergence is
\[
\chi^2(Q_{\mathrm{risk},1}\Vert Q_{\mathrm{risk},0})
:=
\sum_{\substack{o\in O\\Q_{\mathrm{risk},0}\{o\}>0}}
\frac{\bigl(Q_{\mathrm{risk},1}\{o\}
             -Q_{\mathrm{risk},0}\{o\}\bigr)^2}
     {Q_{\mathrm{risk},0}\{o\}}.
\]
In sums over the entire space, a summand with zero denominator is assigned zero; absolute continuity makes its numerator zero as well. Expanding the square and using that both laws have total mass one gives
\[
\begin{aligned}
1+\chi^2(Q_{\mathrm{risk},1}\Vert Q_{\mathrm{risk},0})
&=
\sum_{o\in O}
\frac{Q_{\mathrm{risk},1}\{o\}^2}
     {Q_{\mathrm{risk},0}\{o\}}\\
&=
\frac12+
\left(\frac12-u_{\mathrm{risk}}\right)^2+
\left(\frac12+u_{\mathrm{risk}}\right)^2\\
&=1+2u_{\mathrm{risk}}^2.
\end{aligned}
\]
In particular,
\[
\chi^2(Q_{\mathrm{risk},1}\Vert Q_{\mathrm{risk},0})
=2u_{\mathrm{risk}}^2
\le16u_{\mathrm{risk}}^2.
\]
% lean: CausalSmith.Stat.MarNearcompleteFrontier.parametric_oneCell_pair

\item \emph{Product total variation.}
We use total variation with the convention
\[
\operatorname{TV}
 (Q_{\mathrm{risk},1}^{\otimes n},Q_{\mathrm{risk},0}^{\otimes n})
=
\sup_{\mathcal E\subseteq O^n}
\left|
Q_{\mathrm{risk},1}^{\otimes n}(\mathcal E)
-
Q_{\mathrm{risk},0}^{\otimes n}(\mathcal E)
\right|.
\]
Absolute continuity is preserved by these finite products. Factoring the finite sum of squared product masses divided by reference masses yields
\[
1+\chi^2
 (Q_{\mathrm{risk},1}^{\otimes n}\Vert Q_{\mathrm{risk},0}^{\otimes n})
=
\left\{
1+\chi^2(Q_{\mathrm{risk},1}\Vert Q_{\mathrm{risk},0})
\right\}^{n}.
\]
The inequality \(1+t\le e^t\) for \(t\ge0\), followed by the one-record bound, therefore gives
\[
\begin{aligned}
1+\chi^2
 (Q_{\mathrm{risk},1}^{\otimes n}\Vert Q_{\mathrm{risk},0}^{\otimes n})
&\le
\exp\!\left\{
n\chi^2(Q_{\mathrm{risk},1}\Vert Q_{\mathrm{risk},0})
\right\}\\
&\le \exp(16n u_{\mathrm{risk}}^2),
\end{aligned}
\]
so
\[
\chi^2
 (Q_{\mathrm{risk},1}^{\otimes n}\Vert Q_{\mathrm{risk},0}^{\otimes n})
\le \exp(16n u_{\mathrm{risk}}^2)-1.
\]
On a finite space, total variation equals half the sum of the absolute mass differences. Applying Cauchy--Schwarz to that sum, with the reference masses as weights, gives
\[
\begin{aligned}
\operatorname{TV}
 (Q_{\mathrm{risk},1}^{\otimes n},Q_{\mathrm{risk},0}^{\otimes n})
&\le
\frac12\sqrt{
\chi^2
 (Q_{\mathrm{risk},1}^{\otimes n}\Vert Q_{\mathrm{risk},0}^{\otimes n})
}\\
&\le
\frac12\sqrt{\exp(16n u_{\mathrm{risk}}^2)-1}.
\end{aligned}
\]
Finally, the calibrated budget implies
\[
\exp(16n u_{\mathrm{risk}}^2)\le e^{\log 2}=2,
\qquad
\operatorname{TV}
 (Q_{\mathrm{risk},1}^{\otimes n},Q_{\mathrm{risk},0}^{\otimes n})
\le\frac12.
\]
% lean: CausalSmith.Stat.MarNearcompleteFrontier.parametric_product_tv_at_radius

\item \emph{The squared-error lower bound.}
The estimator class in \cref{obj:synth_6} is nonempty, since it contains
\[
T_{\mathrm{zero}}(o):=0
\qquad(o\in O^n).
\]
Fix \(T\in\mathcal T_n\), and define its two error events by
\[
\mathcal E_{\mathrm{risk},\iota}
:=
\left\{
o\in O^n:
\frac{u_{\mathrm{risk}}}{2}
\le
\left|T(o)-\tau(P_{\mathrm{risk},\iota})\right|
\right\},
\qquad \iota\in\{0,1\}.
\]
The targets are separated by \(u_{\mathrm{risk}}\). The triangle inequality therefore gives
\[
\mathcal E_{\mathrm{risk},0}^{\,c}
\subseteq \mathcal E_{\mathrm{risk},1}.
\]
Using this inclusion and the total-variation bound,
\[
\begin{aligned}
&Q_{\mathrm{risk},0}^{\otimes n}(\mathcal E_{\mathrm{risk},0})
+
Q_{\mathrm{risk},1}^{\otimes n}(\mathcal E_{\mathrm{risk},1})\\
&\qquad\ge
Q_{\mathrm{risk},0}^{\otimes n}(\mathcal E_{\mathrm{risk},0})
+
Q_{\mathrm{risk},0}^{\otimes n}(\mathcal E_{\mathrm{risk},1})
-
\operatorname{TV}
 (Q_{\mathrm{risk},0}^{\otimes n},Q_{\mathrm{risk},1}^{\otimes n})\\
&\qquad\ge\frac12.
\end{aligned}
\]
Consequently,
\[
\max_{\iota\in\{0,1\}}
Q_{\mathrm{risk},\iota}^{\otimes n}(\mathcal E_{\mathrm{risk},\iota})
\ge\frac14.
\]
For each \(\iota\in\{0,1\}\), nonnegativity of squared error gives the threshold bound
\[
\left(\frac{u_{\mathrm{risk}}}{2}\right)^2
Q_{\mathrm{risk},\iota}^{\otimes n}(\mathcal E_{\mathrm{risk},\iota})
\le
\mathbb E_{Q_{\mathrm{risk},\iota}^{\otimes n}}
\left[
\{T-\tau(P_{\mathrm{risk},\iota})\}^2
\right].
\]
If the error probability under the first law is at most that under the second, then
\[
\mathbb E_{Q_{\mathrm{risk},1}^{\otimes n}}
\left[
\{T-\tau(P_{\mathrm{risk},1})\}^2
\right]
\ge
\left(\frac{u_{\mathrm{risk}}}{2}\right)^2\frac14
=
\frac{u_{\mathrm{risk}}^2}{16}.
\]
If that ordering is reversed, the corresponding bound is
\[
\mathbb E_{Q_{\mathrm{risk},0}^{\otimes n}}
\left[
\{T-\tau(P_{\mathrm{risk},0})\}^2
\right]
\ge
\left(\frac{u_{\mathrm{risk}}}{2}\right)^2\frac14
=
\frac{u_{\mathrm{risk}}^2}{16}.
\]
Thus, in either case,
\[
\max_{\iota\in\{0,1\}}
\mathbb E_{Q_{\mathrm{risk},\iota}^{\otimes n}}
\left[
\{T-\tau(P_{\mathrm{risk},\iota})\}^2
\right]
\ge\frac{u_{\mathrm{risk}}^2}{16}.
\]

For every \(P\in\mathcal M(d,q)\), binary potential outcomes imply
\[
-1\le Y(1)-Y(0)\le1,
\qquad
-1\le\tau(P)\le1.
\]
Since \(T\) also takes values in \([-1,1]\),
\[
0\le
\mathbb E_{P_O^{\otimes n}}\{T-\tau(P)\}^2
\le4.
\]
The worst-law supremum therefore exists as a finite real number and dominates the risks at both constructed laws. Taking the estimator infimum in \cref{obj:def:point-risk} gives
\[
\mathfrak R^*_{n,d,q}
\ge\frac{u_{\mathrm{risk}}^2}{16}
=
\frac{a_{\mathrm{risk}}^2}{16n}
=
\frac{c_0}{n}.
\]
% lean: CausalSmith.Stat.MarNearcompleteFrontier.parametric_point_pair_lower

\item \emph{Coverage-dependent calibration.}
Fix \(\alpha\in(0,1/2)\), and define
\[
b_{\mathrm{cov}}:=1-2\alpha,
\qquad
a_{\mathrm{len}}
:=
\min\left\{
\frac14,\frac{\sqrt{\log(1+b_{\mathrm{cov}}^2)}}4
\right\},
\qquad
c_{0,\alpha}:=\frac{a_{\mathrm{len}}b_{\mathrm{cov}}}{2}.
\]
Here
\[
b_{\mathrm{cov}}>0,
\qquad
\log(1+b_{\mathrm{cov}}^2)>0,
\qquad
0<a_{\mathrm{len}}\le\frac14,
\qquad
c_{0,\alpha}>0.
\]
The minimum also gives
\[
a_{\mathrm{len}}^2
\le\frac{\log(1+b_{\mathrm{cov}}^2)}{16}.
\]
Define
\[
u_{\mathrm{len}}:=\frac{a_{\mathrm{len}}}{\sqrt n}.
\]
As before, \(\sqrt n\ge1\) and \((\sqrt n)^2=n\), so
\[
0<u_{\mathrm{len}}\le\frac14,
\qquad
n u_{\mathrm{len}}^2=a_{\mathrm{len}}^2,
\qquad
16n u_{\mathrm{len}}^2\le\log(1+b_{\mathrm{cov}}^2).
\]
Both \(a_{\mathrm{len}}\) and \(c_{0,\alpha}\) depend only on \(\alpha\).
% lean: CausalSmith.Stat.MarNearcompleteFrontier.parametric_radius_budget_at_tolerance

\item \emph{The interval pair and its total variation.}
Using the same one-cell family, define
\[
P_{\mathrm{len},0}:=P_{1/2},
\qquad
P_{\mathrm{len},1}:=P_{1/2+u_{\mathrm{len}}},
\qquad
Q_{\mathrm{len},\iota}:=(P_{\mathrm{len},\iota})_O
\quad(\iota\in\{0,1\}).
\]
The construction and the same three-atom calculation give
\[
\begin{gathered}
P_{\mathrm{len},0},P_{\mathrm{len},1}\in\mathcal M(d,q),\\
\tau(P_{\mathrm{len},0})=\frac12,
\qquad
\tau(P_{\mathrm{len},1})=\frac12+u_{\mathrm{len}},\\
Q_{\mathrm{len},1}\ll Q_{\mathrm{len},0},
\qquad
\chi^2(Q_{\mathrm{len},1}\Vert Q_{\mathrm{len},0})
=2u_{\mathrm{len}}^2\le16u_{\mathrm{len}}^2.
\end{gathered}
\]
The product calculation above applies to this pair. The new budget gives
\[
\exp(16n u_{\mathrm{len}}^2)
\le
\exp\{\log(1+b_{\mathrm{cov}}^2)\}
=
1+b_{\mathrm{cov}}^2.
\]
Using symmetry of total variation and \(b_{\mathrm{cov}}>0\),
\[
\begin{aligned}
\operatorname{TV}
 (Q_{\mathrm{len},0}^{\otimes n},Q_{\mathrm{len},1}^{\otimes n})
&=
\operatorname{TV}
 (Q_{\mathrm{len},1}^{\otimes n},Q_{\mathrm{len},0}^{\otimes n})\\
&\le\frac12\sqrt{\exp(16n u_{\mathrm{len}}^2)-1}\\
&\le\frac12\sqrt{b_{\mathrm{cov}}^2}
=\frac{b_{\mathrm{cov}}}{2}
=\frac{1-2\alpha}{2}.
\end{aligned}
\]
% lean: CausalSmith.Stat.MarNearcompleteFrontier.parametric_product_tv_at_tolerance_radius

\item \emph{Coverage forces expected length.}
Fix \(I\in\mathcal C_\alpha(n,d,q)\), with endpoints as in \cref{obj:def:interval-class}, and define
\[
\mathcal E_{\mathrm{cov},\iota}
:=
\left\{
o\in O^n:
\tau(P_{\mathrm{len},\iota})\in[L(o),U(o)]
\right\},
\qquad \iota\in\{0,1\}.
\]
Honesty gives
\[
Q_{\mathrm{len},0}^{\otimes n}(\mathcal E_{\mathrm{cov},0})
\ge1-\alpha,
\qquad
Q_{\mathrm{len},1}^{\otimes n}(\mathcal E_{\mathrm{cov},1})
\ge1-\alpha.
\]
All these events are measurable because the sample space is finite. The definition of total variation gives
\[
\left|
Q_{\mathrm{len},0}^{\otimes n}(\mathcal E_{\mathrm{cov},1})
-
Q_{\mathrm{len},1}^{\otimes n}(\mathcal E_{\mathrm{cov},1})
\right|
\le\frac{b_{\mathrm{cov}}}{2},
\]
and hence
\[
Q_{\mathrm{len},0}^{\otimes n}(\mathcal E_{\mathrm{cov},1})
\ge1-\alpha-\frac{b_{\mathrm{cov}}}{2}.
\]
Taking complements,
\[
Q_{\mathrm{len},0}^{\otimes n}(\mathcal E_{\mathrm{cov},0}^{\,c})
\le\alpha,
\qquad
Q_{\mathrm{len},0}^{\otimes n}(\mathcal E_{\mathrm{cov},1}^{\,c})
\le\alpha+\frac{b_{\mathrm{cov}}}{2}.
\]
The union bound therefore yields
\[
\begin{aligned}
Q_{\mathrm{len},0}^{\otimes n}
 (\mathcal E_{\mathrm{cov},0}\cap\mathcal E_{\mathrm{cov},1})
&=
1-
Q_{\mathrm{len},0}^{\otimes n}
 (\mathcal E_{\mathrm{cov},0}^{\,c}
  \cup\mathcal E_{\mathrm{cov},1}^{\,c})\\
&\ge1-2\alpha-\frac{b_{\mathrm{cov}}}{2}
=\frac{b_{\mathrm{cov}}}{2}.
\end{aligned}
\]
On the intersection, the interval contains both \(1/2\) and \(1/2+u_{\mathrm{len}}\), so
\[
L(o)\le\frac12,
\qquad
U(o)\ge\frac12+u_{\mathrm{len}},
\qquad
U(o)-L(o)\ge u_{\mathrm{len}}.
\]
Consequently,
\[
Q_{\mathrm{len},0}^{\otimes n}
 (\mathcal E_{\mathrm{cov},0}\cap\mathcal E_{\mathrm{cov},1})
\le
Q_{\mathrm{len},0}^{\otimes n}
 \{o:u_{\mathrm{len}}\le U(o)-L(o)\}.
\]
Since interval length is nonnegative, its expectation satisfies
\[
u_{\mathrm{len}}\,
Q_{\mathrm{len},0}^{\otimes n}
 \{o:u_{\mathrm{len}}\le U(o)-L(o)\}
\le
\mathbb E_{Q_{\mathrm{len},0}^{\otimes n}}[U-L].
\]
Combining these bounds,
\[
\begin{aligned}
\frac{u_{\mathrm{len}}b_{\mathrm{cov}}}{2}
&\le
u_{\mathrm{len}}\,
Q_{\mathrm{len},0}^{\otimes n}
 (\mathcal E_{\mathrm{cov},0}\cap\mathcal E_{\mathrm{cov},1})\\
&\le
u_{\mathrm{len}}\,
Q_{\mathrm{len},0}^{\otimes n}
 \{o:u_{\mathrm{len}}\le U(o)-L(o)\}\\
&\le
\mathbb E_{Q_{\mathrm{len},0}^{\otimes n}}[U-L].
\end{aligned}
\]

The honest interval class is nonempty: the procedure
\[
I_{\mathrm{full}}(o):=[-1,1]
\qquad(o\in O^n)
\]
covers every target because \(\tau(P)\in[-1,1]\). For every admissible interval and law,
\[
0\le U(o)-L(o)\le2,
\qquad
0\le\mathbb E_{P_O^{\otimes n}}[U-L]\le2.
\]
Thus the worst-law expected length is finite and dominates its value at \(P_{\mathrm{len},0}\). Taking the interval infimum in \cref{obj:def:length-risk} gives
\[
\mathfrak L^*_{n,d,q,\alpha}
\ge
\frac{u_{\mathrm{len}}b_{\mathrm{cov}}}{2}
=
\frac{a_{\mathrm{len}}b_{\mathrm{cov}}}{2\sqrt n}
=
\frac{c_{0,\alpha}}{\sqrt n}.
\]
Together with the squared-error bound, this proves both assertions.
% lean: CausalSmith.Stat.MarNearcompleteFrontier.parametric_interval_pair_lower
% lean: CausalSmith.Stat.MarNearcompleteFrontier.parametric_floor
\end{enumerate}
\end{proof}
